%
% TU/e Style Master Thesis template for LaTeX
%
% 2021 Crated by Marko Boon to new corporate identity,
% Based on a template by Thijs Nugteren and Joos Buijs.
%
% THIS IS THE MAIN FILE (i.e. compile this file, compiling the others directly won't work)
%
\documentclass[a4paper,11pt,twoside]{book}
% \documentclass[a4paper]{article}


\usepackage[backref=true,style=alphabetic]{biblatex} 
\addbibresource{references.bib}
\usepackage{tuethesis2021}




\usepackage[
    contentBlocks,
    debugExtensions,
    definitionLists,
    fancy_lists,
    fencedCode,
    hashEnumerators,
    inlineNotes,
    jekyllData,
    expectJekyllData=false,
    notes,
    pipeTables,
    rawAttribute,
    headerAttributes,
    smartEllipses,
    strikeThrough,
    subscripts,
    superscripts,
    tableCaptions,
    taskLists,
    citations,
    hybrid,
    relativeReferences,
    texComments,
    % theme=witiko/graphicx/http,
]{markdown} 



\begingroup
\catcode`\@=11
\catcode`\%=12
\catcode`\^^A=14
\global\def\markdownRendererImagePrototype#1#2#3#4{^^A
{^^A
\everyeof={\noexpand}^^A
\begin{figure}^^A
    \begin{center}
        \includegraphics[width=0.5\textwidth]{#3}^^A
        \caption{#4}^^A
        \label{#1}^^A
    \end{center}
\end{figure}^^A
}^^A
}
\endgroup



\markdownSetup{theme=witiko/graphicx/http}


\def\markdownRendererJekyllDataStart{ }
\def\markdownRendererJekyllDataEnd{ }


%
% These commands need to be defined in order to produce a correct and personalized document.
% You will get error messages if not all these commands are defined. 
% If you don't need a specific command, just leave the argument empty. 
% For example: \docsubtitle{}
%     
\newcommand{\doctitle}{Secure Sessions for Ad Hoc Multiparty Computation in MPyC}
\newcommand{\docsubtitle}{Preparation phase report}
\newcommand{\me}{Emil Nikolov}
\newcommand{\studentNumber}{0972305}
\newcommand{\email}{emil.e.nikolov@gmail.com}
\newcommand{\version}{Draft version}
\newcommand{\placeMonthYear}{Eindhoven, November 2022}
\newcommand{\department}{Department of Mathematics and Computer Science}
\newcommand{\group}{Coding Theory and Cryptology Group}
\newcommand{\firstCommitteeMember}{Dr. ir. L.A.M. (Berry) Schoenmakers} % use all the titles for your committee members!
\newcommand{\secondCommitteeMember}{Dr. Mike Holenderski} % usually the daily supervisor
\newcommand{\thirdCommitteeMember}{Dr. Savio Sciancalepore} % usually the external member
\newcommand{\builddate}{\today} 
\def\markdownOptionOutputDir{build}

\usepackage{mdframed} % surround environments with tags
% \usepackage[linenos=true]{minted} % code snippet colors
\usepackage[outputdir=build]{minted} % code snippet colors
\usemintedstyle{xcode}
% \usemintedstyle{paraiso-light}
% \usemintedstyle{arduino} % code snippet color scheme

\usepackage{csquotes}
\usepackage[english]{babel} % english rules

\usepackage{myglossary} % internal package wrapping glosseries

\begin{document}

\definecolor{LightGray}{gray}{0.9}


\frontmatter

\let\inputmintedorg\inputminted
\renewcommand{\inputminted}[2]{%
    \begin{mdframed}
        \inputmintedorg[linenos=true,fontsize=\footnotesize,breaklines=true,breakanywhere=true]{#1}{#2}
    \end{mdframed}
}


%---------------------------------------TITLE PAGE-----------------------------------%


\begin{titlepage}
    \begin{center}
        \vspace{1.6cm}
        \includegraphics[height=2cm]{figures/TUe-logo-descriptor-line-scarlet-rgb}\\
        \vspace{1.6cm}

        \Large
        \department\\
        \group\\
        \vspace*{1cm}
        \Huge
        \textbf{\doctitle}
        \vspace{0.5cm}

        \LARGE
        \docsubtitle
        \vspace{1.5cm}

        \large
        \textbf{\me}\\
        \vspace{0.4cm}
        Id nr: \studentNumber \\
        \texttt{\email}
        \vfill


        Supervisor : \firstCommitteeMember\\

        \vfill

        Committee members:\\
        \firstCommitteeMember\\
        \secondCommitteeMember\\
        \thirdCommitteeMember\\
        \vfill
        \builddate \\

    \end{center}
\end{titlepage}

% \clearpage

%Sometimes line numbers are nice, uncomment the next line to enable:
%\linenumbers
%---------------------------------------ABSTRACT-----------------------------------%

% \begin{abstract}
\chapter*{Abstract}
\markdownInput{notion/00-abstract.md}
% \end{abstract}
% \markdownInput[slice=^ ^section2]{notion/00-abstract.md}




%---------------------------------------CONTENTS-----------------------------------%

\tableofcontents

%---------------------------------------ABBREVIATIONS-----------------------------------%

\printnoidxglossary[type=\acronymtype,title=List of Abbreviations]
\listoffigures

% \lstlistoflistings
% \listoftables
%---------------------------------------TABLE OF CONTENTS-----------------------------------%


\mainmatter


%---------------------------------------CHAPTERS-----------------------------------%


\markdownInput{notion/01-introduction.md}
% \markdownInput{notion/02-goals.md}

\setkeys{Gin}{scale=0.25}
\markdownInput{notion/03-technical-survey.md}
% \markdownInput{notion/04-design.md}
\markdownInput{notion/05-implementation.md}
\markdownInput{notion/06-conclusions.md}


%---------------------------------------BIBLIOGRAPHY-----------------------------------%

\printbibliography
% \bibliographystyle{plain}
% \bibliographystyle{plainurl}
% \bibliography{references}



%\appendix
%\addcontentsline{toc}{chapter}{Appendix}

\end{document}
